\section{The signed-volume trap}
\label{sec:trap}

\subsection{Where the polynomial stops being a volume}

The functional $D$ of \eqref{eq:signed-volume} is a polynomial on the whole of
$(\mathbb{S}^{d-1})^{n}$, but $d!\,\vol_{d}(P) = D(p)$ only when $p$ realises
$\mathcal{C}$ as a convex polytope. Away from that set the oriented surface
prescribed by $\mathcal{C}$ may intersect itself, and $D(p)/d!$ can then exceed
$\vol_{d}(\conv(p))$, so it is the volume of no polytope on those points. Write
$\Omega_{\mathcal{C}} \subseteq (\mathbb{S}^{d-1})^{n}$ for the set of
configurations that do realise $\mathcal{C}$ convexly, its
\emph{convex-position cell}.

Since $\Omega_{\mathcal{C}}$ is a subset of the domain, maximising $D$ over the
larger set can only give a larger value:
\begin{equation}
  \label{eq:relaxation}
  \sup_{p \in \Omega_{\mathcal{C}}} \vol_{d}
  \;\;\leq\;\;
  \frac{1}{d!}\max_{p \in (\mathbb{S}^{d-1})^{n}} D(p) .
\end{equation}
This is the one direction that is useful: an upper bound on the polynomial
bounds the geometric maximum of that combinatorial type, with no convexity
reasoning at all. The reverse direction fails, as the rest of this section
shows.

\subsection{The heptagonal bipyramid}

Take $d = 3$, $n = 9$, and the combinatorial type of the heptagonal bipyramid,
degree sequence $4^{7}7^{2}$ (label \texttt{plantri-9-24} in our enumeration).
Berman and Hanes showed that a double $m$-pyramid with Property Z is unique up
to congruence, with volume $(m/3)\sin(2\pi/m)$ \cite{berman-hanes-1970}; for
$m = 7$ this is
\begin{equation*}
  \tfrac{7}{3}\sin\tfrac{2\pi}{7} \;=\; 1.824273459092 ,
\end{equation*}
the value every convex-position run below reaches \artref{testtrap}.
Block-coordinate ascent on $D$, however, converges from most starts to fixed
points of signed value
\begin{equation*}
  \tfrac{1}{3!}\,D \;=\; \tfrac{7}{3}\sin\tfrac{4\pi}{7} \;=\; 2.2748318,
\end{equation*}
which is not merely above this value but above the global maximum
$3\sqrt{2\sqrt{3}-3} = 2.043750115900$ over all nine-vertex polytopes
\cite{hoehner-ledford-2026}. No polytope of any type attains it; the
configuration has folded through itself.

The value $\tfrac{7}{3}\sin\tfrac{4\pi}{7}$ is the Berman--Hanes formula
$\tfrac{m}{3}\sin\tfrac{2\pi}{m}$ with $2\pi/m$ replaced by $4\pi/m$, and this
is no coincidence. Label the ring of the double $m$-pyramid type
$\mathcal{C}$ so that its facets join ring vertex $j$ to $j+1$, and for
$1 \leq k < m/2$ let $S_{m,k}$ have apexes $\pm e_{3}$ and ring vertices
$p_{j} = (\cos\tfrac{2\pi kj}{m}, \sin\tfrac{2\pi kj}{m}, 0)$, $0 \leq j < m$,
so that $S_{m,1}$ is the convex configuration. When $\gcd(m,k) = 1$ these are
the points of $S_{m,1}$ relabelled: consecutive ring vertices of $\mathcal{C}$
are now $k$ steps apart on the equator, and the ring traces the star polygon
$\{m/k\}$ instead of the $m$-gon. When $g = \gcd(m,k) > 1$ only $m/g$ of the
ring vertices are distinct, and the ring traces the star polygon
$\{(m/g)/(k/g)\}$ $g$ times. The apexes lie at unit height on either side of
the equatorial plane, so each spans a cone of volume $A/3$ over a planar
region of signed area $A$. The ring, made of $m$ triangles from the centre of
area $\tfrac{1}{2}\sin(2\pi k/m)$ each, has signed area
$(m/2)\sin(2\pi k/m)$, which counts each point of the disc with the ring's
winding number around it, $k$ at the centre \artref{testtrap}. Hence
\begin{equation}
  \label{eq:starbipyramid}
  \tfrac{1}{3!}D(S_{m,k}) \;=\; \frac{m}{3}\sin\frac{2\pi k}{m}.
\end{equation}
For $k \geq 2$ the surface prescribed by $\mathcal{C}$ passes through itself:
$\conv(S_{m,k})$ is a regular bipyramid whose boundary does not use the facets
of $\mathcal{C}$, so $S_{m,k} \notin \Omega_{\mathcal{C}}$. The configuration
reached above is $S_{7,2}$, the bipyramid over the heptagram $\{7/2\}$.

Two features of this matter. First, the fixed points of the ascent step
$p_{i} \mapsto n_{i}/|n_{i}|$ are exactly the configurations with
$p_{i} = n_{i}/|n_{i}|$ for every $i$, and $S_{m,k}$ is one of them
\artref{testtrap}. So every converged point, folded or not, meets the
first-order condition of Lemma~\ref{lem:propertyz} to convergence tolerance;
that condition is necessary for a maximiser but does not distinguish one from
a fold.

Second, the folds attract a large share of starts. Restricting to
convex-position fixed points and recording how many restarts must be discarded
makes this visible across the bipyramid family \artref{testtrap}: with $150$
restarts from a fixed seed,
\begin{center}
\begin{tabular}{lccccc}
  \toprule
  double $m$-pyramid & $m=3$ & $m=4$ & $m=5$ & $m=6$ & $m=7$ \\
  \midrule
  restarts discarded (of 150)      & 0 & 0 & 18 & 81 & 132 \\
  discarded value $\tfrac{1}{3!}D$ & --- & --- & 0.979642 & 1.732051 & 2.274832 \\
  convex value $\tfrac{1}{3!}D$    & 0.866025 & 1.333333 & 1.585094 & 1.732051 & 1.824273 \\
  \bottomrule
\end{tabular}
\end{center}
so that at $m = 7$ only $18$ of $150$ starts survive; each discarded run here
converged to $S_{m,2}$. Note that the fold need not carry the larger value to
attract starts: at $m = 5$ it carries less than the convex configuration and
still takes $18$, and at $m = 6$ it carries the same value and takes $81$. The
same type in our exhaustive $n = 9$ sweep discards $17$ of its $20$ restarts
\artref{cases}.
In every case all surviving runs reproduce the Berman--Hanes closed form to
better than $10^{-14}$.

\subsection{Consequences}

\begin{enumerate}
  \item A sum-of-squares or moment relaxation applied to the bare polynomial
    $D$ cannot help here. Any valid upper bound it certifies for this type is
    at least $2.2748$, hence above the true nine-vertex maximum $2.0438$, and
    so cannot exclude the type. Convexity has to enter the formulation, not be
    checked afterwards.
  \item Every candidate must be certified to lie in $\Omega_{\mathcal{C}}$
    before its value is reported. A type whose ascent yields no convex-position
    fixed point is recorded as infeasible rather than silently assigned a
    folded value.
\end{enumerate}

Inequality \eqref{eq:relaxation} nonetheless survives, and
Section~\ref{sec:triage} turns it into the partition of types on which the
results below are reported.
